\documentclass[11pt,a4paper]{article}

% ------------------------------------------------------------------ Pakete
\usepackage[utf8]{inputenc}
\usepackage[T1]{fontenc}
\usepackage[ngerman]{babel}
\usepackage{lmodern}
\usepackage{amsmath}
\usepackage{booktabs}
\usepackage{array}
\usepackage{graphicx}
\usepackage{caption}
\usepackage{geometry}
\geometry{margin=2.4cm}
\usepackage{xcolor}
\usepackage{titlesec}
\usepackage[hidelinks]{hyperref}
\usepackage{parskip}

\definecolor{secblue}{RGB}{31,73,125}
\titleformat{\section}{\normalfont\large\bfseries\color{secblue}}{\thesection.}{0.6em}{}
\titleformat{\subsection}{\normalfont\normalsize\bfseries}{\thesubsection}{0.6em}{}
\captionsetup{font=small,labelfont=bf}

\begin{document}

% ------------------------------------------------------------------ Titel
\begin{center}
{\LARGE\bfseries Lösung von AsteroidStatic mit Deep Reinforcement Learning}
\end{center}
\vspace{0.2cm}
{\color{gray}
Task 2 \textendash{} Deep Reinforcement Learning (Prof.\ Thomas Nierhoff, OTH Amberg-Weiden)\\
Autor: Himanshu Mohan Nikam\\
Matrikelnummer: 25901848
}
\vspace{0.4cm}

% ================================================================== 1
\section{Projektüberblick}
Das Projekt besteht aus mehreren eigenen Dateien. Die vorgegebene Datei \texttt{asteroid.py}
wurde nicht verändert:
\begin{itemize}\setlength\itemsep{1pt}
  \item \texttt{stoch\_actor\_critic.py} \textendash{} stochastischer Actor-Critic (One-step, on-policy).
  \item \texttt{DDPG.py} \textendash{} Deep Deterministic Policy Gradient.
  \item \texttt{TD3.py} \textendash{} Twin Delayed DDPG.
  \item \texttt{tune\_stochastic\_ac.py}, \texttt{tune\_ddpg.py}, \texttt{tune\_td3.py} \textendash{}
        Optuna-Tuning für die jeweilige Methode.
  \item \texttt{eval\_*.py} sowie die zugehörigen Tuning-Logs, Optuna-Datenbanken
        (\texttt{*.db}), die besten Konfigurationen (\texttt{best\_*.json}) und die
        Auswertungsergebnisse (\texttt{eval\_*\_results.json}).
\end{itemize}

Gemessen werden die beiden geforderten Zielgrößen, jeweils gemittelt über zehn Arenen
(Seeds 0 bis 9): erstens die Anzahl der Umgebungsaufrufe bis zum erstmaligen Erreichen des
Ziels (Metrik 1, Stichprobeneffizienz) und zweitens die kürzeste erfolgreiche Episodenlänge
(Metrik 2, Geschwindigkeit). Zusätzlich wird die Anzahl der Aufrufe festgehalten, bei der
diese kürzeste Episode zuerst auftrat. Für das Tuning werden andere Seeds (115 und 119)
benutzt als für die finale Auswertung.

% ================================================================== 2
\section{Ansätze}
Die Umgebung besitzt einen kontinuierlichen Aktionsraum, weshalb alle drei umgesetzten
Verfahren direkt auf dem Vektor $a=[a_x,a_y]$ arbeiten. Untersucht wurden ein on-policy
Verfahren und zwei off-policy Verfahren.

\subsection{On-policy: Stochastischer Actor-Critic (in stoch\_actor\_critic.py)}
Dieses Verfahren folgt der Vorlesungsformulierung des One-step Actor-Critic. Der Actor ist
eine stochastische Gauß-Politik $\pi(a|s,\theta)$, die Mittelwert und Streuung ausgibt; die
Aktion wird aus dieser Verteilung gezogen. Der Critic schätzt den Zustandswert $V(s,w)$.
Beide werden mit demselben TD-Fehler
$\delta = R + \gamma \hat V(S',w) - \hat V(S,w)$ aktualisiert, der Critic über $\delta^2$ und
der Actor über $-\log \pi(a|s)\,\delta$. Zusätzlich wird ein Entropiebonus addiert, der die
Streuung der Politik offen hält, sowie der Faktor $I \leftarrow \gamma I$ aus der
Vorlesungsfolie verwendet. Die Exploration entsteht hier ausschließlich durch das Ziehen aus
der Verteilung.

\subsection{Off-policy: DDPG (in DDPG.py)}
DDPG ersetzt die stochastische Politik durch einen deterministischen Actor $\mu(s)$ und den
Zustandswert durch einen Aktionswert $Q(s,a)$. Dadurch wird das Verfahren off-policy und kann
einen Replay Buffer nutzen, was für die Stichprobeneffizienz wichtig ist. Der Actor wird über
den deterministischen Policy-Gradienten trainiert, also in Richtung steigender $Q$-Werte.
Stabilisiert wird das Ganze durch Zielnetze mit weichen Polyak-Updates und durch zeitlich
korreliertes Ornstein-Uhlenbeck-Rauschen, das die Exploration liefert.

\subsection{Off-policy: TD3 (in TD3.py)}
TD3 erweitert DDPG um drei Punkte. Erstens werden zwei Critics gelernt und im Ziel das
Minimum beider verwendet, was die typische Überschätzung der Aktionswerte dämpft. Zweitens
wird auf die Zielaktion ein begrenztes Rauschen addiert (Target Policy Smoothing), damit der
Critic keine scharfen Ausreißer im $Q$-Verlauf ausnutzen kann. Drittens wird der Actor
seltener aktualisiert als die Critics, damit die Wertschätzung dem Actor vorauseilt.

\subsection{Weitere erprobte Ansätze}
Vor den drei eingereichten Verfahren wurde auch die DQN-Familie umgesetzt und getestet
(DQN, Double DQN und Dueling DDQN). Da DQN einen endlichen Aktionsraum benötigt, wurde die
Aktion dafür in neun Bang-Bang-Aktionen diskretisiert, also alle Kombinationen aus
$\{-1,0,+1\}$ für $a_x$ und $a_y$. Dieser Ansatz löste einzelne Arenen zuverlässig, ist aber
nicht Teil der Abgabe, da der Schwerpunkt auf den Verfahren für kontinuierliche Aktionen
liegt.

% ================================================================== 3
\section{Wichtige Techniken}
\begin{itemize}\setlength\itemsep{3pt}
  \item \textbf{Eingangsnormalisierung.} Positionen und Geschwindigkeiten werden anhand der
        bekannten Grenzen der Umgebung auf einen vergleichbaren Bereich skaliert. Es handelt
        sich um eine reine Skalierung der Eingänge und nicht um eine Analyse der Umgebung. Ohne
        diesen Schritt war keines der Verfahren stabil trainierbar.
  \item \textbf{Asymmetrisches Reward-Clipping.} Die Kollisionsstrafe von $-100$ ist rund zwei
        Größenordnungen größer als der Fortschrittsterm von $0{,}5 \cdot$ Fortschritt. Dadurch
        kollabieren die gelernten Werte auf einen nahezu konstanten negativen Wert, und der
        Agent kann gute von schlechten Aktionen nicht mehr unterscheiden. Die Strafe wird
        deshalb nach unten begrenzt, während der Zielbonus von $+10$ unverändert groß bleibt.
        So bleiben die seltenen erfolgreichen Übergänge das dominierende Signal.
  \item \textbf{Exploration.} Ein rein zufälliger Agent erreicht das Ziel nur in etwa
        $0{,}2\,\%$ der Episoden und auf mehreren Arenen überhaupt nicht. Die Exploration muss
        daher dauerhaft stark bleiben. Beim stochastischen Actor-Critic übernimmt das der
        Entropiebonus, bei DDPG und TD3 das korrelierte Ornstein-Uhlenbeck-Rauschen.
  \item \textbf{Begrenzung der Streuung.} Die Streuung der Gauß-Politik wird über
        \texttt{log\_std} gelernt und auf einen festen Bereich begrenzt. Ohne diese Begrenzung
        wächst sie durch den Entropiebonus unbegrenzt an, bis die Verteilung ungültig wird
        (siehe Abschnitt 6).
  \item \textbf{Stabilität.} Zusätzlich werden Gradienten begrenzt, Zielnetze verwendet und der
        Replay Buffer vor dem ersten Update mit zufälligen Übergängen gefüllt.
\end{itemize}

% ================================================================== 4
\section{Hyperparameter Tuning}
Die Hyperparameter werden mit Optuna gesucht, für jede Methode in einem eigenen Skript und
einer eigenen SQLite-Datenbank, sodass unterbrochene Läufe fortgesetzt werden können. Getunt
werden unter anderem die Lernraten von Actor und Critic, $\gamma$, die Netzgröße sowie die
methodenspezifischen Parameter, also $\tau$ und die Rauschstärke bei DDPG und TD3
beziehungsweise der Entropiekoeffizient beim stochastischen Actor-Critic.

Die Tuning-Seeds sind von den Auswertungs-Seeds getrennt. Um geeignete Tuning-Seeds zu finden,
wurden die Kandidaten 100 bis 119 zunächst mit einem Zufallsagenten vermessen. Auf den meisten
davon wird das Ziel nie zufällig erreicht, sodass das Tuning dort kein verwertbares Signal
liefert. Gewählt wurden deshalb die Seeds 115 und 119, auf denen der Zufallsagent das Ziel
zumindest gelegentlich erreicht. Die Zielfunktion minimiert die Aufrufe bis zur ersten Lösung;
Konfigurationen, die innerhalb des Episodenbudgets nicht lösen, werden ersatzweise nach ihrem
geringsten Abstand zum Ziel bewertet, damit die Suche auch dann eine Richtung hat.

\begin{table}[h]
\centering
\small
\caption{Beste von Optuna gefundene Konfiguration je Verfahren.}
\label{tab:hyper}
\begin{tabular}{lccc}
\toprule
Hyperparameter & Stoch.\ Actor-Critic & DDPG & TD3 \\
\midrule
Actor-Lernrate       & offen & $1{,}7\!\cdot\!10^{-5}$ & $1{,}4\!\cdot\!10^{-4}$ \\
Critic-Lernrate      & offen & $4{,}8\!\cdot\!10^{-4}$ & $1{,}9\!\cdot\!10^{-3}$ \\
$\gamma$             & offen & $0{,}992$ & $0{,}962$ \\
$\tau$ (Polyak)      & --    & $1{,}8\!\cdot\!10^{-3}$ & $8{,}0\!\cdot\!10^{-3}$ \\
Rauschstärke         & offen & $\sigma_{\text{OU}}=0{,}37$ & $\sigma_{\text{OU}}=0{,}32$ \\
Target Policy Smoothing & -- & -- & $0{,}34$ / $0{,}50$ \\
Policy Delay         & --    & -- & $1$ \\
Batch-Größe          & offen & $128$ & $128$ \\
Netzgröße            & offen & $256$ & $256$ \\
\bottomrule
\end{tabular}
\end{table}

% ================================================================== 5
\section{Ergebnisse}
Jedes Verfahren wird mit seiner besten Konfiguration auf den Seeds 0 bis 9 ausgewertet, mit
einem Budget von 5000 Episoden je Seed. Tabelle~\ref{tab:main} zeigt, wie viele Arenen gelöst
wurden und die Mittelwerte über die jeweils gelösten Seeds. Die Auswertung des stochastischen
Actor-Critic steht noch aus und wird nachgetragen.

\begin{table}[h]
\centering
\small
\caption{Ergebnisse auf den Auswertungs-Seeds 0 bis 9, gemittelt über die jeweils gelösten Seeds.}
\label{tab:main}
\begin{tabular}{lcccc}
\toprule
Verfahren & Gelöste Seeds & Metrik 1 & Metrik 2 & Zugehörige Aufrufe \\
\midrule
Stoch.\ Actor-Critic & offen          & offen  & offen  & offen \\
DDPG                 & 5/10           & 4.473  & 24{,}2 & 88.761 \\
TD3                  & \textbf{10/10} & 89.295 & 37{,}9 & 147.919 \\
\bottomrule
\end{tabular}
\end{table}

Die Mittelwerte beziehen sich jeweils nur auf die von einem Verfahren gelösten Seeds und sind
deshalb nicht unmittelbar vergleichbar. TD3 löst als einziges Verfahren auch die beiden
schwersten Arenen (Seeds 4 und 8), die lange Wege von 83 beziehungsweise 92 Schritten
erfordern; genau diese beiden Läufe heben die Mittelwerte von TD3 deutlich an. DDPG löst nur
fünf, eher leichte Arenen, was seinen niedrigen Wert bei Metrik 1 erklärt. Tabelle
\ref{tab:perseed} zeigt daher die Einzelergebnisse.

\begin{table}[h]
\centering
\small
\caption{Ergebnisse je Seed. Ein Strich bedeutet, dass die Arena im Budget nicht gelöst wurde.}
\label{tab:perseed}
\begin{tabular}{ccccc}
\toprule
Seed & \multicolumn{2}{c}{DDPG} & \multicolumn{2}{c}{TD3} \\
     & Metrik 1 & Metrik 2 & Metrik 1 & Metrik 2 \\
\midrule
0 & --    & -- & 121.118 & 26 \\
1 & --    & -- & 32.362  & 23 \\
2 & 1.235 & 27 & 10.022  & 25 \\
3 & --    & -- & 250.901 & 35 \\
4 & --    & -- & 144.963 & 83 \\
5 & 3.185 & 26 & 3.192   & 25 \\
6 & --    & -- & 45.808  & 24 \\
7 & 7.228 & 23 & 3.151   & 23 \\
8 & 4.573 & 22 & 279.128 & 92 \\
9 & 6.146 & 23 & 2.301   & 23 \\
\bottomrule
\end{tabular}
\end{table}

Betrachtet man nur die fünf Seeds, die beide Verfahren lösen (2, 5, 7, 8 und 9), so liegt DDPG
bei Metrik 1 im Mittel bei 4.473 Aufrufen und TD3 bei 59.559. Ohne den teuren Seed 8 sinkt TD3
allerdings auf 4.667 Aufrufe und liegt damit gleichauf. Bei Metrik 2 sind beide Verfahren auf
diesen Seeds nahezu identisch.

% Platzhalter für Abbildungen, nach dem Erzeugen der Plots einkommentieren:
% \begin{figure}[h]\centering
%   \includegraphics[width=0.8\linewidth]{lernkurve.png}
%   \caption{Lernkurve: Episodenlänge und Erfolgsrate über der Anzahl der Aufrufe.}
% \end{figure}

% ================================================================== 6
\section{Fehlgeschlagene Versuche}

\subsection{Deterministischer Actor-Critic}
Zuerst wurde ein deterministischer Actor-Critic umgesetzt, also im Kern DDPG ohne dessen
Stabilisierungen. Dieser kollabierte zuverlässig auf die Nullaktion: Der Actor gab eine
Beschleunigung nahe null aus, das Raumschiff blieb am Startpunkt stehen und überlebte bis zum
Zeitlimit mit einer Belohnung von etwa null. Das ist aus Sicht des Agenten rational, weil
Nichtstun keine Kollisionsstrafe kostet, während jeder Versuch, das Ziel zu erreichen, im
Hindernisfeld endet. Erst eine deutlich verstärkte Exploration brachte ihn aus diesem lokalen
Optimum heraus. Für die Abgabe wurde er durch den stochastischen Actor-Critic ersetzt, der
näher an der Vorlesungsformulierung liegt.

\subsection{Symmetrisches Reward-Clipping}
Ein erster Versuch begrenzte die Belohnung symmetrisch auf $[-1,1]$. Damit wurde zwar die
Kollisionsstrafe gedämpft, gleichzeitig aber auch der Zielbonus auf denselben Betrag
gestaucht. Da das Ziel nur äußerst selten erreicht wird, verlor es dadurch jede Sonderstellung
gegenüber einer Kollision, und die Wertschätzungen kollabierten weiterhin. Erst die
asymmetrische Variante aus Abschnitt 3 löste das Problem.

\subsection{Zu weite Suchbereiche beim Tuning}
Ein Tuning-Lauf des stochastischen Actor-Critic brach nach zwei Trials vollständig ab. Die
Suchbereiche für die Lernraten waren zu weit gewählt, sodass Optuna Werte bis 0,245 für den
Critic vorschlug. Bei solchen Lernraten divergiert der TD-Fehler, und der Entropiebonus treibt
die gelernte Streuung der Politik unbegrenzt nach oben, bis sie numerisch überläuft und die
Normalverteilung ungültig wird. Da Optuna Ausnahmen standardmäßig weiterreicht, beendete
dieser eine Trial die gesamte Studie. Abhilfe schaffen engere Suchbereiche, die Begrenzung von
\texttt{log\_std} und das Abfangen fehlgeschlagener Trials, damit die Suche weiterläuft.

\subsection{Schwache Exploration}
Kleines Gauß-Rauschen um eine bereits kollabierte deterministische Politik reicht nicht aus.
Der Agent zittert dann nur lokal um den Startpunkt und legt nie genug Strecke zurück, um das
Ziel überhaupt zu finden. Notwendig war entweder korreliertes Rauschen oder das Einstreuen von
Zufallsaktionen mit voller Magnitude.

% ================================================================== 7
\section{Diskussion}
Die größte Schwierigkeit dieser Aufgabe liegt in der Exploration und nicht in der Lernregel.
Das Ziel liegt hinter einem dichten Hindernisfeld, und der Fortschrittsterm der Belohnung zeigt
geradlinig darauf zu. Zustände näher am Ziel tragen dadurch ein höheres Kollisionsrisiko und
einen niedrigeren Wert, sodass ein Agent es rational vermeidet, sich dem Ziel zu nähern. Genau
das erzeugt die beiden beobachteten lokalen Optima, nämlich die Flucht in den freien Raum und
das Verharren am Startpunkt.

Unter den ausgewerteten Verfahren ist TD3 mit Abstand das robusteste und löst als einziges alle
zehn Arenen. Die drei Erweiterungen gegenüber DDPG wirken hier also genau wie beabsichtigt: Das
Minimum zweier Critics verhindert, dass der Actor überschätzte Aktionswerte ausnutzt, was gerade
bei der großen Kollisionsstrafe entscheidend ist. DDPG ist auf den Arenen, die es löst, sehr
stichprobeneffizient, scheitert aber an den schwereren Layouts. Bemerkenswert ist, dass die
getunte TD3-Konfiguration einen Policy Delay von eins gewählt hat, die verzögerte Aktualisierung
des Actors auf den Tuning-Seeds also keinen Vorteil brachte.

Erwartungsgemäß dürfte der stochastische Actor-Critic bei Metrik 1 schlechter abschneiden als
die beiden off-policy Verfahren, da er jeden Übergang nur einmal verwendet und keinen Replay
Buffer besitzt. Auf einer Aufgabe, bei der das Ziel so selten erreicht wird, wiegt das schwer,
weil eine seltene erfolgreiche Episode sofort wieder verworfen wird, statt vielfach
wiederverwendet zu werden.

% ================================================================== 8
\section{Fazit}
Verglichen wurden ein on-policy Verfahren und zwei off-policy Verfahren für kontinuierliche
Aktionen, zusammen mit den Techniken, die das Lernen in dieser Umgebung überhaupt ermöglichen,
also Eingangsnormalisierung, asymmetrisches Reward-Clipping und dauerhaft starke Exploration.
TD3 ist das robusteste Verfahren und löst alle zehn Arenen, während DDPG nur die leichteren
Layouts bewältigt, dort aber sehr stichprobeneffizient ist. Die Zahlen für den stochastischen
Actor-Critic werden nach Abschluss seiner Auswertung ergänzt.

% ================================================================== 9
\section{Quellen}

\end{document}
